\documentclass[12pt, a4paper]{article}
\usepackage[utf8]{inputenc}
\usepackage[T1]{fontenc}
\usepackage{amsmath, amssymb, amsthm}
\usepackage{geometry}
\geometry{margin=2.5cm}
\usepackage{hyperref}
\usepackage{booktabs}
\usepackage{xeCJK}
\setCJKmainfont{cwTeX Q Ming}

\title{矩陣微分 Cheat Sheet}
\author{Sheng-Yuan Chen}
\date{2026-04-07}

\begin{document}
\maketitle

\section*{前言}

在研讀計量經濟學，特別是\textbf{因子模型（Factor Models）}與\textbf{主成分分析（PCA）}時，矩陣微分是一個不可缺少的工具。本文整理常用公式，統一採用\textbf{分子佈局（Numerator Layout）}：對行向量求導結果為行向量，對矩陣求導結果與分子同形。

\section{跡（Trace）的微分規則}

跡的兩個核心性質：
\begin{itemize}
  \item \textbf{轉置不變性：} $\text{tr}(A) = \text{tr}(A^T)$
  \item \textbf{循環性質（Cyclic Property）：} $\text{tr}(ABC) = \text{tr}(BCA) = \text{tr}(CAB)$
\end{itemize}

\subsection*{常用公式}

$$\frac{\partial \, \text{tr}(AX)}{\partial X} = A^T$$

$$\frac{\partial \, \text{tr}(X^T A)}{\partial X} = A$$

$$\frac{\partial \, \text{tr}(AX^T B)}{\partial X} = B^T A^T$$

\section{二次型（Quadratic Form）的微分}

\subsection*{向量形式}

$$\frac{\partial (x^T A x)}{\partial x} = (A + A^T)x$$

若 $A$ 對稱（如共變異數矩陣），則化簡為 $2Ax$。

\subsection*{矩陣跡形式}

$$\frac{\partial \, \text{tr}(X^T A X)}{\partial X} = (A + A^T)X$$

若 $A$ 對稱（如 $X^T X$），化簡為 $2AX$。

\section{行列式與逆矩陣（Determinant \& Inverse）}

\subsection*{對數行列式}

$$\frac{\partial \ln |X|}{\partial X} = (X^{-1})^T$$

若 $X$ 對稱（如協方差矩陣），化簡為 $X^{-1}$。

\subsection*{逆矩陣的跡}

$$\frac{\partial \, \text{tr}(A X^{-1})}{\partial X} = -(X^{-1} A X^{-1})^T$$

\section{實際應用：因子模型（Factor Models）}

PCA 的最佳化問題為：

$$\max_{F} \; \text{tr}\!\left(F^T (X^T X) F\right) \quad \text{subject to} \quad F^T F = I$$

令 $A = X^T X$（對稱矩陣），對 $F$ 微分並加入 Lagrange 乘數項，整理得一階條件：

$$AF = FV$$

這正是\textbf{特徵值問題（Eigenvalue Problem）}：$F$ 的各行是 $A$ 的特徵向量，$V$ 是對角的特徵值矩陣。

\section{小結}

\begin{table}[h]
\centering
\begin{tabular}{lll}
\toprule
類型 & 公式 & 備註 \\
\midrule
線性跡 & $\frac{\partial \, \text{tr}(AX)}{\partial X} = A^T$ & 基本型 \\
二次型（向量） & $\frac{\partial x^T A x}{\partial x} = (A+A^T)x$ & 對稱時為 $2Ax$ \\
二次型（矩陣跡） & $\frac{\partial \, \text{tr}(X^T AX)}{\partial X} = (A+A^T)X$ & 對稱時為 $2AX$ \\
對數行列式 & $\frac{\partial \ln |X|}{\partial X} = (X^{-1})^T$ & 對稱時為 $X^{-1}$ \\
逆矩陣跡 & $\frac{\partial \, \text{tr}(AX^{-1})}{\partial X} = -(X^{-1}AX^{-1})^T$ & MLE 常見 \\
PCA 一階條件 & $AF = FV$ & 特徵值問題 \\
\bottomrule
\end{tabular}
\end{table}

\end{document}
